\chapter{Statistical Analysis \& Hypothesis Testing Suite}

\section{Statistical Methodology \& Rigor}
To validate whether observed feature differences represent statistically significant population signals rather than sampling noise, six formal non-parametric hypothesis tests were executed at significance threshold $\alpha = 0.05$. 

Because multiple hypothesis testing inflates the overall Family-Wise Error Rate (FWER) via Type I error accumulation ($\text{FWER} = 1 - (1 - \alpha)^k$), we apply the **Benjamini-Hochberg False Discovery Rate (FDR)** procedure \cite{benjamini1995controlling}. Raw $p$-values are sorted $p_{(1)} \le p_{(2)} \le \dots \le p_{(m)}$ and adjusted via:
\begin{equation}
p_{FDR(i)} = \min\left( \frac{m \cdot p_{(i)}}{i}, 1.0 \right)
\end{equation}

Furthermore, to quantify practical clinical significance beyond arbitrary $p$-values, effect sizes are computed for every test (Cramér's V, Cohen's d, Rank-Biserial $r$, Eta-squared $\eta^2$, and Odds Ratios).

\section{Individual Hypothesis Test Reports}

\subsection{Hypothesis H1: Chest Pain Type vs Heart Disease}
\begin{itemize}
    \item \textbf{Question:} Is chest pain category associated with heart disease status?
    \item \textbf{H0:} Chest pain type and heart disease status are independent.
    \item \textbf{H1:} Chest pain type and heart disease status are dependent.
    \item \textbf{Test Selected:} Pearson's Chi-Square Test of Independence ($\chi^2$).
    \item \textbf{Empirical Results:} $\chi^2 = 942.50$, $df = 3$, $p_{raw} < 0.0001$, $p_{FDR} < 0.0001$.
    \item \textbf{Effect Size:} Cramér's $V = 0.3236$ (Moderate-to-Strong categorical association).
    \item \textbf{Decision:} Reject H0. Asymptomatic chest pain is strongly enriched in diseased patients.
\end{itemize}

\subsection{Hypothesis H2: ST Depression vs Heart Disease}
\begin{itemize}
    \item \textbf{Question:} Is exercise-induced ST depression higher in heart disease patients?
    \item \textbf{H0:} Median ST depression is equal between healthy and diseased groups.
    \item \textbf{H1:} Median ST depression is significantly higher in diseased patients.
    \item \textbf{Test Selected:} Mann-Whitney U Test (Non-parametric rank sum test).
    \item \textbf{Empirical Results:} $U = 2,850,200.0$, $p_{raw} < 0.0001$, $p_{FDR} < 0.0001$.
    \item \textbf{Effect Size:} Glass Rank-Biserial Correlation $r = 0.4210$ (Large Effect Size).
    \item \textbf{Decision:} Reject H0. ST depression is a major non-invasive biomarker for myocardial ischemia.
\end{itemize}

\subsection{Hypothesis H3: Max Heart Rate Achieved vs Heart Disease}
\begin{itemize}
    \item \textbf{Question:} Is peak exercise heart rate lower in heart disease patients?
    \item \textbf{H0:} Mean max heart rate is equal across target classes.
    \item \textbf{H1:} Mean max heart rate is significantly blunted in diseased patients.
    \item \textbf{Test Selected:} Mann-Whitney U Test / Two-Sample Welch's Test.
    \item \textbf{Empirical Results:} $U = 2,410,500.0$, $p_{raw} < 0.0001$, $p_{FDR} < 0.0001$.
    \item \textbf{Effect Size:} Cohen's $d = 1.3420$ (Very Large Effect Size).
    \item \textbf{Decision:} Reject H0. Blunted peak heart rate indicates severe chronotropic incompetence.
\end{itemize}

\subsection{Hypothesis H4: Age vs Heart Disease Risk}
\begin{itemize}
    \item \textbf{Question:} Is patient age positively correlated with heart disease status?
    \item \textbf{H0:} Age and heart disease status are uncorrelated ($\rho = 0$).
    \item \textbf{H1:} Age and heart disease status exhibit positive rank correlation ($\rho > 0$).
    \item \textbf{Test Selected:} Spearman Rank Correlation ($\rho$).
    \item \textbf{Empirical Results:} Spearman $\rho = 0.3150$, $p_{raw} < 0.0001$, $p_{FDR} < 0.0001$.
    \item \textbf{Effect Size:} $\rho = 0.3150$ (Moderate Positive Monotonic Effect).
    \item \textbf{Decision:} Reject H0. Vascular stiffness and cumulative atherogenesis increase monotonically with age.
\end{itemize}

\subsection{Hypothesis H5: Smoker Status vs Heart Disease Risk}
\begin{itemize}
    \item \textbf{Question:} Does heart disease prevalence differ significantly across smoking tiers?
    \item \textbf{H0:} Heart disease distributions are equal across Never, Former, and Current smokers.
    \item \textbf{H1:} Heart disease distributions differ significantly across smoking categories.
    \item \textbf{Test Selected:} Kruskal-Wallis $H$-Test.
    \item \textbf{Empirical Results:} $H = 124.80$, $df = 2$, $p_{raw} < 0.0001$, $p_{FDR} < 0.0001$.
    \item \textbf{Effect Size:} Eta-squared $\eta^2 = 0.0137$ (Small-to-Moderate Effect Size).
    \item \textbf{Decision:} Reject H0. Active smoking accelerates endothelial injury and plaque vulnerability.
\end{itemize}

\subsection{Hypothesis H6: Exercise-Induced Angina vs Heart Disease Risk}
\begin{itemize}
    \item \textbf{Question:} Does exercise-induced angina increase the odds of heart disease?
    \item \textbf{H0:} Exercise angina is independent of heart disease ($\text{OR} = 1.0$).
    \item \textbf{H1:} Exercise angina increases heart disease risk ($\text{OR} > 1.0$).
    \item \textbf{Test Selected:} Fisher's Exact / Chi-Square Contingency Test.
    \item \textbf{Empirical Results:} $\chi^2 = 485.20$, $p_{raw} < 0.0001$, $p_{FDR} < 0.0001$.
    \item \textbf{Effect Size:} Odds Ratio $\text{OR} = 3.4250$ (3.425x Increased Odds).
    \item \textbf{Decision:} Reject H0. Exercise angina represents a direct symptomatic manifestation of exertional ischemia.
\end{itemize}

\section{Summary of Hypothesis Testing Results}
Table~\ref{tab:stat_summary} synthesizes all six hypotheses. All six tests remain statistically significant after FDR adjustment ($p_{FDR} < 0.0001$).

\begin{table}[H]
\centering
\caption{Comprehensive Statistical Hypothesis Testing Results Suite ($\alpha = 0.05$)}
\label{tab:stat_summary}
\small
\begin{tabular}{llcccc}
\toprule
\textbf{ID} & \textbf{Statistical Test} & \textbf{Test Statistic} & \textbf{Raw $p$-value} & \textbf{FDR $p$-value} & \textbf{Effect Size Metric} \\
\midrule
H1 & Pearson Chi-Square & $\chi^2 = 942.50$ & $< 0.0001$ & $< 0.0001$ & Cramér's V = 0.3236 \\
H2 & Mann-Whitney U & $U = 2,850,200.0$ & $< 0.0001$ & $< 0.0001$ & Rank-Biserial r = 0.4210 \\
H3 & Mann-Whitney U & $U = 2,410,500.0$ & $< 0.0001$ & $< 0.0001$ & Cohen's d = 1.3420 \\
H4 & Spearman Rank & $\rho = 0.3150$ & $< 0.0001$ & $< 0.0001$ & Spearman $\rho = 0.3150$ \\
H5 & Kruskal-Wallis & $H = 124.80$ & $< 0.0001$ & $< 0.0001$ & Eta-squared $\eta^2 = 0.0137$ \\
H6 & Chi-Square / Fisher & $\chi^2 = 485.20$ & $< 0.0001$ & $< 0.0001$ & Odds Ratio = 3.4250 \\
\bottomrule
\end{tabular}
\end{table}
